\section{Prompt-conditional distribution matching with MMD}
\label{sec:mmd}
\noindent\textbf{Evidence status.} This section gives the original
conditional-MMD method and preparation checks. Completed policy results,
calibrated-kernel follow-ups, and diagnostics appear in
Section~\ref{sec:newresults}. This section distinguishes the
implemented two-reference target from the earlier proposal to generate
larger reference sets with an existing frozen teacher. JSD is documented
as a comparison candidate, not an executed conditional training arm.

\subsection{Separating reference choice from discrepancy choice}
A single word PMF pooled across unrelated questions may reward vocabulary
that is inappropriate for an individual answer. Replacing KL with a
bounded or kernel discrepancy does not repair that target mismatch.
The conditional proposal therefore makes two separate choices:
\begin{enumerate}
\item Define a reference distribution for the exact prompt or conversation.
\item Compare generated answers with that reference using a categorical
discrepancy such as JSD, or a feature-space discrepancy such as MMD.
\end{enumerate}
For example, answers to ``What is a Jupyter notebook?'' should be compared
with valid answers to that question, not with pooled vocabulary from
unrelated programming tasks. Features may capture lexical frequencies,
code/prose proportions, formatting, or semantic embeddings. Their choice
defines which properties are controlled; these targets are not
interchangeable notions of successful sampling.

\subsection{Implemented reference subset: two acceptable answers per question}
The current implementation avoids new teacher generation. Starting from
the prepared programming split (228 training, 16 validation, 32 test
prompts), it retains questions with exactly two distinct source answers
whose human coherence, correctness, and helpfulness ratings are each
at least 3 on the 0--4 scale. These answers are model-written and
human-rated, not newly collected human-written demonstrations.
The filtered counts are shown in Table~\ref{tab:mmddata}.

\begin{table}[htbp]
\centering\small
\caption{Verified conditional-MMD preparation statistics. Reference answers
are cached source examples; none is a new student or teacher generation.}
\label{tab:mmddata}
\begin{tabular}{@{}lrr@{}}
\toprule
Split & Exact questions & Reference answers\\
\midrule
Training & 102 & 204\\
Validation & 10 & 20\\
Test & 19 & 38\\
\midrule
Total & 131 & 262\\
\bottomrule
\end{tabular}
\end{table}

Splits inherit the programming dataset's conversation-family separation.
The implementation additionally keys every comparison by its exact
canonical conversation ID, rather than by topic or a shared opening.
Training rejects held-out reference IDs; validation/test references are
used only to evaluate their corresponding questions. The original
programming word PMF remains a training-only diagnostic and is not the
semantic optimization target.

Let $a_{c,1},a_{c,2}$ denote the two acceptable answers and $e(a)$ their
frozen embeddings. The actual target is
\begin{equation}
 P_c^{\mathrm{ref}}=\tfrac12\delta_{e(a_{c,1})}
                        +\tfrac12\delta_{e(a_{c,2})}.
 \label{eq:mmdtarget}
\end{equation}
Equal weights are an experimental choice. Two acceptable answers do not
estimate the frequencies of all valid human answer styles. This target
therefore represents two prototypes, not a recovered human distribution.

\subsection{Frozen full-answer features and the configured kernel}
The encoder is the pinned
\texttt{sentence-transformers/all-MiniLM-L6-v2}, which provides
384-dimensional representations \cite{minilmcard}. Encoder weights stay
frozen. To avoid dropping long-answer tails, tokenization is performed
without truncation, followed by nonoverlapping chunks fitting within
256 encoder tokens including special tokens. Each chunk has
attention-masked mean pooling and unit normalization. Chunk vectors
are combined using their content-token counts as weights, then
renormalized:
\begin{equation}
 u_j=\frac{\operatorname{MeanPool}(h_j)}
              {\|\operatorname{MeanPool}(h_j)\|_2},\qquad
 e(a)=\frac{\sum_j L_j u_j}{\|\sum_j L_j u_j\|_2},
 \label{eq:chunkembedding}
\end{equation}
where $L_j$ is the chunk's content length, with weight one for an empty
chunk. This preserves all visible token chunks, but averaging can lose
global order and dependencies. It is not a proof that code semantics
or correctness are preserved.

The kernel is a mixture of three Gaussian RBF kernels:
\begin{equation}
 k(u,v)=\frac13\sum_{s\in\{0.5,1,2\}}
      \exp\!\left(-\frac{\|u-v\|_2^2}{2b^2s^2}\right).
 \label{eq:mmdkernel}
\end{equation}
The fixed base bandwidth squared is the median nonzero within-question
squared distance between the two training-reference embeddings only:
$b^2=0.2786315084$. No validation/test pair sets this bandwidth.
The median training-pair cosine similarity is 0.8606843. These values
describe the cached reference representation, not model-generation
quality. Both the encoder and reference vectors remain fixed in RL;
new student answers still require encoder inference.

\subsection{Population MMD and the implemented finite estimator}
For independent reference features $X,X'\sim P_c$ and generated
features $Z,Z'\sim Q_{\theta,c}$, the squared MMD is
\begin{align}
 \operatorname{MMD}_k^2(P_c,Q_{\theta,c})
 &=\|\mu_{P_c}-\mu_{Q_{\theta,c}}\|_{\mathcal H_k}^2\nonumber\\
 &=\E k(X,X')+\E k(Z,Z')-2\E k(X,Z),
 \label{eq:mmdpopulation}
\end{align}
where $\mu_P=\E_{X\sim P}k(X,\cdot)$ is the kernel mean embedding
\cite{gretton2012}. The two within-distribution terms and the cross term
compare both spread and cross-group similarity. For suitable
characteristic kernels, zero population MMD identifies equality of
feature distributions. A noninjective text encoder can still map
different or incorrect answers to similar features, so equality in
feature space does not imply equality of text distributions or correctness.

For $n$ generated features $z_1,\ldots,z_n$ and $m=2$ fixed reference
features $x_1,x_2$, the implementation uses
\begin{equation}
 \widehat D_U=
 \frac{\sum_{i\ne j}k(z_i,z_j)}{n(n-1)}
 +\frac{\sum_{a,b=1}^{m}k(x_a,x_b)}{m^2}
 -\frac{2\sum_{i=1}^{n}\sum_{a=1}^{m}k(z_i,x_a)}{nm}.
 \label{eq:mmdestimator}
\end{equation}
Only the generated self-term is a U-statistic; the reference term
includes diagonal terms because it defines the fixed empirical target
in Equation~\eqref{eq:mmdtarget}. Conditional on that target and
independent generated answers, $\E\widehat D_U=
\operatorname{MMD}_k^2(P_c^{\mathrm{ref}},Q_{\theta,c})$.
This is not an unbiased estimate against an unknown population of human
answers. The finite estimate can be negative and is not clipped to zero.
Evaluation also records a nonnegative V-statistic version that replaces
the first denominator by $n^2$ and includes generated diagonal terms.
The two estimators must not be silently interchanged.

MMD does not require estimating a full categorical word PMF and avoids
the specific logarithmic singularity of unsmoothed forward KL. It still
depends on the representation, kernel, bandwidth, reference coverage,
and generated sample count. Minimizing feature MMD does not imply that
the previous full-word KL objective will improve.

\subsection{Conditional policy objective and leave-one-answer-out credit}
The configured batch has $G=4$ distinct questions and $n=4$ answers per
question, giving 16 generated answers across two GPU ranks. Every
semantic comparison and quality baseline remains within its exact
question. The actor starts from the pinned Qwen2.5-3B-Instruct checkpoint
with a fresh LoRA adapter. It reuses the frozen
\texttt{Skywork-Reward-V2-Qwen3-8B} scorer and its existing audit/calibration;
it does not retrain the 3B Bradley--Terry scorer described in the earlier
protocol. This is an independently prepared profile, not a continuation
of the numerical SFT adapter.

The updated quality term preserves the score of a capped answer and
subtracts a separate nontermination penalty:
\begin{equation}
 r_Q(c,y)=\sigma\!\left(\frac{r_\phi(c,y)-m_\phi}{s_\phi}\right)
             -0.1\,\ind\{y\text{ reaches its generation cap}\}.
 \label{eq:mmdquality}
\end{equation}
It therefore differs from the earlier protocol's EOS-multiplied reward.
The standardized score is clipped to $[-30,30]$ before the sigmoid.
The objective combines average quality, average conditional MMD, and
the separate sequence-policy regularizer:
\begin{equation}
 J=\E\!\left[\frac1{Gn}\sum_{g,i}r_Q(c_g,Y_{g,i})
          -\frac{\lambda}{G}\sum_g\widehat D_{U,g}\right]
       -\beta\E_c\KL(\pi_\theta(\cdot\mid c)\Vert\pi_0(\cdot\mid c)),
 \label{eq:mmdobjective}
\end{equation}
with configured $\lambda=0.5$ and $\beta=0.002$.
The MMD weight is an initial setting: a coefficient of 0.5 has a different
scale from 0.5 times word KL and does not establish the same quality trade-off.

Writing $D_g$ for the full-question estimate and $D_{g,-i}$ for the
estimate after removing answer $i$ and recomputing all generated
normalizers, the credits are
\begin{align}
 A^D_{g,i}&=n(D_{g,-i}-D_g),\\
 A^Q_{g,i}&=r_Q(c_g,Y_{g,i})-
               \frac1{n-1}\sum_{j\ne i}r_Q(c_g,Y_{g,j}),\\
 \widehat{\mathcal L}_{\mathrm{MMD}}
 &=-\frac1{Gn}\sum_{g,i}\sg\!\left[A^Q_{g,i}+\lambda A^D_{g,i}
       -\beta\log\frac{\pi_\theta(Y_{g,i}\mid c_g)}{\pi_0(Y_{g,i}\mid c_g)}\right]
       \log\pi_\theta(Y_{g,i}\mid c_g).
 \label{eq:mmdsurrogate}
\end{align}
The generated U-statistic requires at least two answers, so its
leave-one-out credit requires at least three. Four answers satisfy both
requirements. Scores and embeddings are treated as detached reward
inputs; gradients flow through actor completion log probabilities,
not through sampled text or the frozen encoder. Global DDP aggregation
collects the groups before computing credits. The configured run uses
up to 200 updates, learning rate $5\times10^{-6}$, 20 warmup updates,
one-passage gradient microbatches, temperature one without top-$p$/top-$k$
truncation, and at most 1,024 new actor tokens. The original run actually
completed 94 updates and 1,504 training answers. Validation/test use
four answers per question, giving 40/76 generated answers per policy.

\subsection{What the current CPU checks establish}
The recorded CPU audit compares cached training-reference embeddings
under three constructions: two copies of each prototype; four copies
of one prototype; and four copies of a reference from another question.
The balanced construction has lower estimated MMD than either
comparison for all 102 training questions. This is a controlled
representation/kernel check using reference-derived points. It is not
an evaluation of generated answers, independently annotated semantic
quality, or learned stochasticity. In particular, preferring balanced
prototype copies over collapsed copies follows naturally from the
kernel objective and does not validate code correctness.
Seven CPU tests passed in the inspected snapshot, covering exact expected
MMD policy gradients, conditional credit, full-text chunk coverage,
reference separation, and two-rank aggregation. They do not supply a
trained-policy result.

Reported prototype coverage assigns generated embeddings to their nearest
reference by cosine similarity and records how many of the two prototypes
are represented. This measures coverage of those prototypes, not the
number or frequencies of true semantic modes. Held-out references remain
unavailable to training. Policy quality, conditional MMD, independent
rubric scores, generation length, termination, and word-distribution
diagnostics must all be examined before a policy benefit can be claimed.
The configured single-policy profile has no matched quality-only control
or second seed. It reserves a 75-minute policy stage within a two-H200,
two-hour allocation template. These are configured budget limits; observed
process times and achieved updates are in Section~\ref{sec:newresults}.
Cross-run differences from
the unigram profile would also change the target, references, scorer,
and batching, so they could not isolate MMD's causal contribution.

\subsection{Frozen-teacher references and conditional JSD: proposed comparisons}
The original discussion proposed a larger prompt-specific reference
set sampled from an existing frozen teacher. This would not train a
new teacher. References could be generated once, quality-checked,
cached, and reused after unloading the teacher. Twenty prompts with
eight answers each require 160 reference generations; 228 prompts
with eight or 16 answers require 1,824 or 3,648. These are illustrative
counts, not performed runs or GPU-time estimates. The implemented
two-human-rated-reference version instead requires zero new teacher
generations, but inherits its narrow prototype coverage.

For categorical prompt-specific word laws, the proposed JSD comparison is
\begin{equation}
 M_c=\tfrac12(P_c+Q_c),\qquad
 \operatorname{JSD}(P_c,Q_c)=
 \tfrac12\KL(P_c\Vert M_c)+\tfrac12\KL(Q_c\Vert M_c).
 \label{eq:conditionaljsd}
\end{equation}
With natural logarithms it lies in $[0,\log 2]$ and is finite despite
missing categories, using the usual zero-mass convention \cite{lin1991}.
A common category partition, including a category for novel words,
must still be specified. A bounded discrepancy does not resolve
reference mismatch or finite-sample noise. Lexical JSD and embedding
MMD control different properties, so their numerical scales and meanings
are not directly interchangeable.

The proposed sequence is to check agreement between independent teacher
reference batches on the same prompts, then compare conditional JSD
and MMD on identical references and student draws. Any subsequent
training comparison should include SFT on those same reference answers,
plus a quality-only control and matched training/evaluation budgets.
Code execution or task-specific correctness checks are useful when
available for programming answers. Teacher bias, missing valid answer
styles, frozen-encoder errors, tiny per-question batches, one seed, and
public-data overlap remain limitations. Neither MMD nor JSD supplies
a trustworthy target by itself, replaces the quality reward, or
guarantees conditional randomness. Larger frozen-teacher reference sets,
conditional-JSD training, and matched SFT/quality-only controls remain
unrun. Post hoc lexical JSD audits in Section~\ref{sec:newresults} are
not a JSD policy-training experiment.
